\section{Model and Identification}\label{sec:model_identification}

% Our theoretical framework is closely related to the two-period \cite{royThoughtsDistributionEarnings1951} models in \cite{lagakosSelectionAgricultureCrossCountry2013} and \cite{lagakosMigrationCostsObservational2020}. 
% In \cite{lagakosSelectionAgricultureCrossCountry2013}, cross-country productivity gaps in agriculture stem from the presence of a subsistence consumption requirement for the agricultural good and heterogeneous worker productivity across the agricultural and non-agricultural sector.
% In \cite{lagakosMigrationCostsObservational2020}, workers with different migration costs and returns sort themselves across multiple regions, which generates different aggregate estimates for rural-urban consumption gaps. 
% We build on these models to develop a multi-period migration model that allows for two-way migration between rural and urban areas, which matches the migration patterns in our data. 
% We maintain the assumption that workers are endowed with location-specific skills, but we introduce different assumptions on the role that preferences and prices play.
% \textcolor{red}{Should we be more explicit here about what these modifications are?}
% The resulting migration decision rule in our setting therefore maps onto a correlated random coefficient model in the vein of \cite{lemieuxEstimatingEffectsUnions1998} and \cite{suriSelectionComparativeAdvantage2011}. 

We study an economy where individuals choose between working in a rural or urban labor market, based on their skills and preferences. 
As is common in the literature, we motivate our empirical strategy with a generalized \cite{royThoughtsDistributionEarnings1951} model.\footnote{See for example \citealt{bazziSkillTransferabilityMigration2016, alvarezAgriculturalWageGap2020, lagakosSelectionAgricultureCrossCountry2013,pulidoBarriersMobilitySorting2021}.} 
In the classical Roy model, people sort across sectors solely based on income or consumption.
Our generalized Roy model adopts a standard generalization that allows location choices to also depend on non-pecuniary factors. 
In the context of migration, the non-pecuniary component of utility captures proximity to family, local amenities, or other idiosyncratic factors that affect location choices but are not reflected in consumption.

\subsection{Consumption and Location Choice}\label{location-choice-model}

The model has three main components.
We first describe how consumption depends on worker productivity across rural and urban locations.
We then decompose productivity into absolute and comparative advantage to characterize sorting and returns to migration.
Finally, we introduce utility shocks and derive the worker’s location choice rule.


\subsubsection{Consumption and Productivity Across Sectors}

The rural and urban labor markets, indexed by $l=R,U$, demand different skills and compensate these skills differently. 
For worker $i$ at time $t$, who works in location $l$, let $y_{it}^l$ denote the systematic component of log consumption.
We model this object as
\begin{align}
    y_{it}^{U}&=\beta_t^U+x_{it}'\gamma^U+\theta_i^U \label{eq:y_it_U}\\
    y_{it}^{R}&=\beta_t^R+x_{it}'\gamma^R+\theta_i^R \label{eq:y_it_R}
\end{align}

\noindent where $x_{i t}$ is a vector of observable characteristics, $\beta^{l}_{t}$ captures the average log consumption in sector $l$ at time $t$, and $\theta^{l}_{i}$ is a time-invariant unobservable that captures worker $i$'s productivity in sector $l$.
The difference between these two sector-specific unobservables, $\theta_{i}^U-\theta_{i}^R$, represents person $i$'s comparative advantage in favor of urban relative to rural work.

\subsubsection{Comparative Advantage and Productivity Decomposition}

Following \cite{lemieuxEstimatingEffectsUnions1998} and \cite{suriSelectionComparativeAdvantage2011}, and building on the correlated random effects framework of \citet{chamberlain1984panel}, we project the time-invariant, sector-specific unobservables onto comparative advantage.
We define $b_U$ and $b_R$ as coefficients from the linear projection of $\theta^{l}_{i}$ onto $\theta^{U}_{i} - \theta^{R}_{i}$, assuming zero means.  
Specifically, we write
% which allows us to simplify the model by focusing on the worker's comparative advantage, and yields the following expressions:
\begin{align}
    \theta_i^U&=b_U(\theta_i^U-\theta_i^R)+\tau_i\\
    \theta_i^R&=b_R(\theta_i^U-\theta_i^R)+\tau_i
\end{align}
\noindent where $\tau_{i}$ captures absolute advantage that affects a worker's productivity in both locations.
This decomposition isolates the part of productivity that drives sorting across locations from the part that raises consumption regardless of where a person works.
By construction, $\tau_{i}$ is orthogonal to comparative advantage.

We can interpret the projection coefficients $b_{U}$ and $b_{R}$ as measuring how strongly a worker's comparative advantage, $\theta_{i}^U-\theta_{i}^R$, maps onto productivity in the urban and rural sectors, respectively.
These coefficients depend on the dispersion of productivity within each sector and on the covariance between urban and rural productivity.
When urban and rural productivities are highly correlated, workers who perform well in one location also tend to perform well in the other, so comparative advantage plays a smaller role in shaping relative productivity.
When the correlation is low, productivity becomes more location-specific, and comparative advantage matters more for sorting across locations.
Similarly, when productivity varies more in one sector than in the other, that sector places greater weight on comparative advantage, making outcomes there more sensitive to differences in workers' skills.

Formally, 
\begin{align}
b_U &= \frac{\sigma_U^2-\sigma_{UR}}{\sigma_U^2+\sigma_R^2-2\sigma_{UR}}, \\
b_R &= \frac{\sigma_{UR}-\sigma_R^2}{\sigma_U^2+\sigma_R^2-2\sigma_{UR}},
\end{align}
where $\sigma_l^2=\text{Var}(\theta_i^l)$ and $\sigma_{UR}=\text{Cov}(\theta_i^U,\theta_i^R)$.

To simplify notation, we define
\[
\theta_i \equiv b_R(\theta_i^U-\theta_i^R),
\]
which is a rescaled measure of comparative advantage that captures both the worker's relative productivity across locations and the extent to which the rural sector rewards that advantage.

We also define
\[
\phi \equiv \frac{b_U-b_R}{b_R},
\]
which measures how much more important comparative advantage is for urban productivity than for rural productivity in the economy.\footnote{The definition requires $b_R\neq 0$, which we assume throughout.}
The parameter $\phi$ is larger when urban productivity varies more across workers or when skills transfer poorly between rural and urban locations, so that relative skill differences matter more in the urban sector.
Conversely, $\phi$ is smaller when productivity varies more in the rural sector or when urban and rural productivities are highly correlated, in which case comparative advantage affects outcomes more symmetrically across locations.
As a result, $\phi$ governs how strongly returns to urban location increase with workers' comparative advantage.

Using these definitions, we can rewrite the systematic components of log consumption as
\begin{align}
y_{it}^{U} &= \beta_t^U + x_{it}'\gamma^U + (1+\phi)\theta_i + \tau_i, \label{eq:ysys_U_phi}\\
y_{it}^{R} &= \beta_t^R + x_{it}'\gamma^R + \theta_i + \tau_i. \label{eq:ysys_R_phi}
\end{align}

The term $\tau_i$ shifts consumption levels equally across locations, while $\theta_i$ governs how relative returns differ between urban and rural work.
When $\phi>0$, comparative advantage enters more strongly in the urban sector, so the returns to urban location rise with $\theta_i$.
In that case, workers with lower rural consumption gain relatively less from migrating.
When $\phi<0$, comparative advantage matters more in rural locations, and the gains from urban migration are larger for workers with lower rural consumption.

\subsubsection{Utility and Location Choice}

We allow for factors other than consumption to affect migration decisions by modeling location choice as a function of utility. 
In addition to consumption, utility depends on non-monetary factors like proximity to family and local amenities.
We capture these with a time- and location-specific idiosyncratic shock, $\nu_{it}^{l}$. 
A person's utility from working in location $l$ at time $t$ is therefore \(V_{it}^{l} = y_{it}^l + \nu_{it}^{l}\). 
These shocks generate variation in location choices even when a person's observable characteristics and productivity remain unchanged.


We assume that $\nu_{it}^l$ is independently and identically distributed across individuals, time periods, and locations. 
This assumption ensures that utility shocks generate transitory movements across locations without introducing persistence or correlation with past choices.
It also implies that $\nu_{it}^{l}$ is orthogonal to the unobserved components of consumption, which simplifies the mapping from the model to the empirical specification.
Under this assumption, migration decisions respond to temporary, idiosyncratic factors that do not reflect systematic differences in productivity across sectors.

An important consequence of this assumption is that the period-by-period location choice rule derived below is equivalent to the solution of a fully forward-looking dynamic optimization problem.
Two conditions deliver this equivalence: the i.i.d.\ assumption ensures that future utility shocks are independent of the current location choice, so there is no informational option value; and the time-invariant productivity assumption (equations~\ref{eq:y_it_U}--\ref{eq:y_it_R}) ensures that payoffs do not depend on past locations, so there is no state dependence in returns.
Together, these eliminate any channel through which today's choice affects future welfare, reducing the dynamic problem to a sequence of static comparisons.
If non-pecuniary factors were instead persistent---or if migration experience altered productivity---the interpretation of trajectories as reflecting comparative advantage would be less precise.

Workers choose their location for period $t$ by comparing expected utility across locations. 
At the end of each period, a worker observes current consumption and utility shocks and forms expectations over future shocks. 
She chooses the urban location in period $t$ if her expected utility in the urban sector exceeds expected utility in the rural sector. 
Following \cite{lemieuxEstimatingEffectsUnions1998}, we assume that the average rural-urban consumption gap,
\begin{equation}
    \beta \equiv \beta_t^U - \beta_t^R,\label{eq:beta}
\end{equation}
is constant over time, while allowing the level of average consumption in each location to vary over time.
Abstracting away from covariates, or equivalently assuming that observable characteristics affect consumption symmetrically across locations, we can rewrite the urban choice rule as\footnote{Covariates $x_{it}$ do not appear in the decision rule because under the symmetric covariate assumption ($\gamma^U = \gamma^R$), their contribution cancels in the urban-rural comparison.} 
\begin{equation}\label{eq:decision-rule}
    E\!\left[V_{it}^U\right] > E\!\left[V_{it}^R\right]
    \;\Longleftrightarrow\;
    \beta + \phi\,\theta_i + E\!\left[\nu_{it}^U - \nu_{it}^R\right] > 0.
\end{equation}
Because the idiosyncratic shocks are i.i.d.\ and productivity is time-invariant, the forward-looking comparison in \eqref{eq:decision-rule} reduces to a sequence of static, period-by-period choices evaluated at the realized shock.
Worker $i$ selects the urban location in period $t$ whenever $\beta + \phi\theta_i + (\nu_{it}^U - \nu_{it}^R) > 0$.
It is this realized variation in $\nu_{it}^U - \nu_{it}^R$, not its expectation, that generates the location switching from which the returns are identified.

\subsection{Empirical Model}\label{subsec:empirical-model}

Let $D_{i t}=1$ indicate that worker $i$ works in the urban location at time $t$, and $D_{i t}=0$ otherwise.
Observed log consumption equals the systematic component described in the model plus an idiosyncratic error term

\begin{equation}
    y_{i t} = y_{i t}^l + \varepsilon_{i t},
\end{equation}
where $l \in \{R,U\}$ denotes location and $\varepsilon_{i t}$ captures transitory shocks to consumption.

Because $D_{it}$ is binary, we can write observed consumption as a random coefficient model in which the return to urban location varies across workers.
In particular, combining the sector-specific consumption equations yields
\begin{align}
    y_{it}&=\beta^R+\theta_i + \tau_i +(\beta+\phi \theta_i)D_{it}+  x_{it}'\gamma^R + D_{it}x_{it}'(\gamma^U-\gamma^R) + \varepsilon_{it},\label{eq:generalized-consumption-eq}
\end{align}
\noindent where $\beta^R$ suppresses the time subscript from equations \eqref{eq:y_it_U}--\eqref{eq:y_it_R} because in estimation, period effects absorb the time-varying level of average consumption.

The worker-specific return to urban location is therefore $\Delta_i \equiv \beta + \phi\theta_{i}$, which depends on comparative advantage.
Since $\theta_{i}$ also affects location choices through the decision rule derived above, the return to urban location is correlated with $D_{it}$.
This places the model in the class of correlated random coefficient (CRC) models.

% Similar to Lagakos and Waugh (2013), Lagakos \emph{et al.} (2020), and Pulido and Swiecki (2021), we assume that unobserved ability is time-invariant.

% This is more general than a standard fixed-effect model since we allow for the possibility that $\theta_{i}$ has a different impact on urban and rural consumption. 

We assume that observable characteristics affect consumption symmetrically across locations, so that $\gamma^U=\gamma^R$.
Under this restriction, equation \eqref{eq:generalized-consumption-eq} simplifies to: 
\begin{align}
    y_{it}&=\beta^R+\theta_i + \tau_i +(\beta+\phi \theta_i)D_{it}+  x_{it}'\gamma + \varepsilon_{it}.\tag{10$'$}\label{eq:generalized-consumption-eq-simple}
\end{align}

Equation \eqref{eq:generalized-consumption-eq-simple} makes clear that heterogeneity in the returns to urban location arises from comparative advantage. 
Workers with different values of $\theta_{i}$ face different gains from working in the urban sector, and these gains are correlated with observed migration choices. 
This highlights a central identification challenge.
Because returns to urban location vary across workers and are correlated with location choices, average returns cannot be recovered from comparisons of urban and rural outcomes without additional structure.
In particular, without further restrictions, returns are only identified for workers who switch locations over time, and counterfactual returns for non-migrants remain unidentified.

In the remainder of this section, we show how this challenge can be addressed by aggregating workers into migration histories and imposing economically motivated restrictions that link returns across subpopulations.


\subsubsection{Unrestricted Group Random Coefficient Model}\label{subsec:unrestricted-grc-model}
The correlated random coefficient (CRC) model in equation \eqref{eq:generalized-consumption-eq-simple} allows worker-specific returns to urban location that are correlated with location choices. 
To take this model to the data, we follow \cite{tjernstromCommentSuri2011} and aggregate the CRC model over discrete migration histories, which yields a  group random coefficient (GRC) representation.
In what follows, we suppress covariates and period effects for clarity; these re-enter the estimation as described above. 

In the GRC framework, we group workers by their observed sequence of location choices over time, which we refer to as migration trajectories.
Different trajectories reflect systematic differences in workers' ability or comparative advantage, to the extent that location choices respond to relative productivity and utility.
As a result, observed migration histories provide a natural way to partition workers into subpopulations with different average outcomes and returns.
Each trajectory has its own average consumption level and return to urban location.
The unrestricted specification places no structure on how these trajectory-specific average returns vary across migration histories.

Because the location indicator $D_{it}$ is binary and the panel has finite length, our data contain a finite number of migration trajectories.
We can summarize each worker's observed sequence of location choices by $\underline d_i\equiv(D_{i1},\dots,D_{iT})$.
These trajectories take values in the finite set $\mathcal D=\{0,1\}^T$, which we use to index subpopulations of workers with the same migration history.

We distinguish between switcher trajectories,
$\mathcal D_S = \left\{ \underline d \in \mathcal D : 0 < \sum_{t=1}^T D_{it} < T \right\},$
and non-switcher trajectories,
$\mathcal D_{NS} = \mathcal D \setminus \mathcal D_S.$
The latter consist of the always-urban trajectory
$d_T = \left\{ \underline d \in \mathcal D : \sum_{t=1}^T D_{it} = T \right\},$ and the always-rural trajectory
$d_N = \left\{ \underline d \in \mathcal D : \sum_{t=1}^T D_{it} = 0 \right\}.$

Without additional restrictions, the CRC model implies the following unrestricted GRC representation for any $t \ge 2$.\footnote{The restriction to $t \ge 2$ reflects the use of the first period to establish baseline trajectory membership. In estimation, we include all available periods and absorb initial-period effects through period indicators.}
\begin{equation}\label{eq:unrestricted-grc-model}
    y_{it}= \sum_{\underline d \in \mathcal D \setminus \{d_T\}}
    \mu_{\underline d}\,\mathbbm{1}\{ \underline d_i=\underline d \} + \sum_{\underline d \in \mathcal D_S}
    \Delta_{\underline d}\, D_{it}\,\mathbbm{1}\{ \underline d_i=\underline d \} + \kappa_{d_T}\, D_{it}\,\mathbbm{1}\{ \underline d_i=d_T \} + \varepsilon_{it}.
\end{equation}

Here, $\mu_{\underline d}$ denotes average rural consumption for workers following trajectory $\underline d$.
In this unrestricted model, average returns to urban location, $\Delta_{\underline d}$, are non-parametrically identified only for switcher trajectories.
For non-switchers, we cannot identify counterfactual outcomes. 
For always-rural workers, the data identify only average rural consumption, $\mu_{d_N}$.
For always-urban workers, we cannot separately identify $\mu_{d_T}$ and $\Delta_{d_T}$ and instead only recover average urban consumption, which we denote by $\kappa_{d_T} \equiv \mu_{d_T} + \Delta_{d_T}$.
This limitation motivates the additional structure introduced in the next subsection.

\subsubsection{Restricted Group Random Coefficient Model}\label{subsec:restricted-grc-model}

To identify returns to urban location for non-migrants, we impose a linear comparative advantage (LCA) restriction that links heterogeneity in migration returns to differences in workers' comparative advantage.
This restriction follows \cite{suriSelectionComparativeAdvantage2011} and builds on the linear correlated random coefficient framework in \cite{lemieuxEstimatingEffectsUnions1998}.

The LCA restriction assumes that the return to urban location varies linearly with a worker's comparative advantage:
\begin{equation}\label{eq:LCA-restriction}
\Delta_i = \beta + \phi\,\theta_i .
\end{equation}

Since comparative advantage also determines average rural consumption in the model, the LCA restriction implies a linear relationship between average returns and average rural consumption across migration trajectories.
As a result, variation in outcomes among switchers identifies the relationship between returns and comparative advantage, which in turn allows us to extrapolate returns for non-migrants.
Since $\tau_i$ does not affect location choices, trajectory-specific differences in $\mu_{\underline{d}}$ identify differences in average comparative advantage.

We do not restrict the sign of this relationship; instead, we allow the data to determine whether returns to migration increase or decrease with comparative advantage.
Our estimation strategy follows \cite{tjernstromCommentSuri2011}, who show how to implement this restriction within a GRC framework using GMM.
For a chosen baseline switcher trajectory $\underline d_0 \in \mathcal D_S$, we can write the restricted GRC model as:
\begin{align}
    y_{it} =& \sum_{\underline d \in \mathcal D \setminus \{d_T\}}
    \mu_{\underline d}\,\mathbbm{1}\{\underline d_i=\underline d\} + \Delta_{\underline d_0} D_{it} \nonumber \\
    & + \sum_{\underline d \in \mathcal D_S \setminus \{\underline d_0\}}
    \phi(\mu_{\underline d}-\mu_{\underline d_0})\,D_{it}\,\mathbbm{1}\{\underline d_i=\underline d\} \nonumber\\
    & + \Big(\mu_{d_T}+\phi(\mu_{d_T}-\mu_{\underline d_0})\Big)\,D_{it}\,\mathbbm{1}\{\underline d_i=d_T\} + \varepsilon_{it}.
    \label{eq:restricted-GRC}
\end{align} 
\noindent for some baseline trajectory $d_{0} \in \mathcal{D}_{S}$.\footnote{The always-urban group ($\underline d_i = d_T$) receives a separate composite coefficient in the last line because these workers have $\sum_t D_{it} = T$ and are not in $\mathcal D_S$.
Their trajectory-specific return is therefore parameterized directly rather than through the switcher summation.}

The relationship in the LCA restriction forms the basis of how we identify returns to non-migrants in our model.
Taking expectations of equation \eqref{eq:LCA-restriction} within migration trajectories implies that for any two switcher trajectories $\underline{d} \neq \underline{d}'$ with distinct average rural consumption, average returns satisfy:

\begin{equation}
    \phi = \frac{\Delta_{\underline{d}} - \Delta_{\underline{d}'}}{\mu_{\underline{d}} - \mu_{\underline{d}'}},\ \ \mu_{\underline{d}} \neq \mu_{\underline{d}'}.\label{eq:phi-proposition1}
\end{equation}

If the LCA parameter, $\phi$, is positive, then urban location has higher returns for the people with higher average consumption in the rural market.
Conversely, a negative $\phi$ would suggest that those who have the lowest rural consumption would receive the highest incremental consumption gains from migrating to the urban market.
We can refer to these two scenarios as migration being either ``pro-rich'' (positive $\phi$) or ``pro-poor'' (negative $\phi$).
 
% Note that a positive \(\phi\) implies a positive correlation between the unobserved location-specific skills in the population, but the converse is not true.
% \footnote{From equation \eqref{eq:phi-variance-covariance}, we have

%   \[\phi = \frac{{Var}\left( \theta_{i}^{U} - \theta_{i}^{R} \right)}{{Cov}\left( \theta_{i}^{U},\theta_{i}^{R} \right) - {Var}\left( \theta_{i}^{R} \right)} > 0 \Leftrightarrow {Cov}\left( \theta_{i}^{U},\theta_{i}^{R} \right) > {Var}\left( \theta_{i}^{R} \right) \Leftrightarrow \rho\  > \frac{\sigma_{R}}{\sigma_{U}},\ \ \rho = \frac{{Cov}\left( \theta_{i}^{U},\theta_{i}^{R} \right)}{\sigma_{U}\sigma_{R}}.\]

%   A positive \(\phi\) requires that (\emph{i}) \(\rho > 0\) since
%   \(\frac{\sigma_{R}}{\sigma_{U}} > 0\), and (\emph{ii})
%   \(\sigma_{U} \geq \sigma_{R}\) since \(|\rho| \leq 1\). That is, we
%   must have a positive correlation between the two location-specific
%   skill sets, and a larger dispersion of urban-specific skills than
%   rural-specific skills. The condition \(\phi < 0\), on the other hand,
%   can occur with \(\rho = 0,\  < 0\) or \(> 0\) and
%   \(\frac{\sigma_{R}}{\sigma_{U}} < 1\) or \(> 1\). In other words,
%   observing \(\phi < 0\) does not allow us to infer the direction of the
%   correlation between location-specific skills, \(\rho\).}


Before turning to estimation, we collect below the assumptions that underlie identification:
\begin{enumerate}[label=\textup{(A\arabic*)},leftmargin=*,itemsep=2pt]
    \item Location-specific productivities $\theta_i^U$ and $\theta_i^R$ are time-invariant.
    \item Non-pecuniary utility shocks $\nu_{it}^l$ are i.i.d.\ across individuals, time, and locations, so that migration trajectories reflect comparative advantage rather than shock persistence.
    \item The average rural-urban consumption gap $\beta$ is constant over time.
    \item Observable characteristics affect consumption symmetrically across locations ($\gamma^U = \gamma^R$).
    \item The linear comparative advantage restriction $\Delta_i = \beta + \phi\,\theta_i$.
\end{enumerate}
% We assess the empirical support for A5 using Hansen's $J$-test in Section \ref{sec:grc-returns}.

We estimate the restricted model in Eq. \eqref{eq:restricted-GRC} using two-step efficient Generalized Method of Moments (GMM), as in \cite{tjernstromCommentSuri2011}. 
Let $\varepsilon_{it}$ denote the residual from equation \eqref{eq:restricted-GRC}. The moment conditions are
\begin{equation}\label{eq:gmm-moments}
    E\!\left[\varepsilon_{it} \cdot z_{it}\right] = 0,
\end{equation}
where the instrument vector $z_{it}$ consists of trajectory indicators $\mathbbm{1}\{\underline{d}_i = \underline{d}\}$, the urban indicator $D_{it}$, interactions $D_{it} \cdot \mathbbm{1}\{\underline{d}_i = \underline{d}\}$ for each switcher trajectory, and the covariate vector $x_{it}$.
The number of moment conditions exceeds the number of parameters, and we test the overidentifying restrictions using Hansen's $J$-statistic.
Standard errors are clustered at the individual level.

While the key identifying variation comes from the switchers in the balanced panel, we also include observations that appear in a subset of panel waves in our estimation.
We include an ``unbalanced'' dummy variable that takes a value of one for individuals with at least one missing round of data, as well as the interaction between the ``unbalanced'' dummy and the urban indicator.\footnote{Since unbalanced individuals lack complete trajectories, all trajectory indicators equal zero for them. They therefore do not contribute to identifying the trajectory-specific parameters $\mu_{\underline{d}}$ or the slope $\phi$. They only help estimate the covariate coefficients $\gamma$ and the unbalanced level effects. This inclusion is valid under two assumptions: (i) attrition is independent of $\varepsilon_{it}$ conditional on $D_{it}$ and $x_{it}$, and (ii) covariate effects on consumption are common across balanced and unbalanced individuals.}
Appendix \ref{app:balanced_panel} shows results using the balanced panel.

